% from ega0-1-fr.tex
\oldpage[0\textsubscript{I}]{44}
canonique (3.7.2)
$\Gamma(V,\mathcal{B})\to\Gamma(U,\psi^*(\mathcal{B}))$, c'est donc aussi un
$\Gamma(V,\mathcal{B})$-module).
\end{sloppypar}
On voit aussitôt que ces homomorphismes sont
compatibles avec les restrictions, et par suite définissent un homomorphisme
canonique fonctoriel
\[
  \gamma:\mathcal{H}\!om_{\mathcal{B}}(\mathcal{M},\mathcal{N})
  \longrightarrow
  \Psi_*\bigl(\mathcal{H}\!om_{\mathcal{A}}
  (\Psi^*(\mathcal{M}),\Psi^*(\mathcal{N}))\bigr)~;
\]
il correspond par suite aussi à cet homomorphisme l'homomorphisme
\[
  \gamma^\sharp:
  \Psi^*\bigl(\mathcal{H}\!om_{\mathcal{B}}(\mathcal{M},\mathcal{N})\bigr)
  \longrightarrow
  \mathcal{H}\!om_{\mathcal{A}}
  (\Psi^*(\mathcal{M}),\Psi^*(\mathcal{N}))
\]
et ces homomorphismes canoniques sont fonctoriels en $\mathcal{M}$ et
$\mathcal{N}$.
\end{env}

\begin{env}[4.4.7]
\label{0.4.4.7-fr}
Supposons que $\mathcal{F}$ (resp. $\mathcal{G}$) soit une
$\mathcal{A}$-Algèbre (resp. une $\mathcal{B}$-Algèbre). Si
$u:\mathcal{G}\to\Psi_*(\mathcal{F})$ est un homomorphisme de
$\mathcal{B}$-Algèbres, $u^\sharp$ est un homomorphisme
$\Psi^*(\mathcal{G})\to\mathcal{F}$ de $\mathcal{A}$-Algèbres~; cela résulte
de la commutativité du diagramme
\[
\begin{tikzcd}[column sep=huge,row sep=large]
\mathcal{G}\otimes_{\mathcal{B}}\mathcal{G}
  \arrow[r] \arrow[d]
  & \mathcal{G} \arrow[d,"u"] \\
\Psi_*(\mathcal{F}\otimes_{\mathcal{A}}\mathcal{F})
  \arrow[r]
  & \Psi_*(\mathcal{F})
\end{tikzcd}
\]
et de (4.4.4). De même, si $v:\Psi^*(\mathcal{G})\to\mathcal{F}$ est un
homomorphisme de $\mathcal{A}$-Algèbres,
$v^\flat:\mathcal{G}\to\Psi_*(\mathcal{F})$ est un homomorphisme de
$\mathcal{B}$-Algèbres.
\end{env}

\begin{env}[4.4.8]
\label{0.4.4.8-fr}
Soient $(Z,\mathcal{C})$ un troisième espace annelé,
$\Psi'=(\psi',\theta')$ un morphisme $(Y,\mathcal{B})\to(Z,\mathcal{C})$, et
$\Psi'':(X,\mathcal{A})\to(Z,\mathcal{C})$ le morphisme composé
$\Psi'\circ\Psi$. Soit $\mathcal{H}$ un $\mathcal{C}$-Module, $u'$ un
homomorphisme de $\mathcal{H}$ dans $\mathcal{G}$~; le composé
$v''=v\circ v'$ est par définition l'homomorphisme de $\mathcal{H}$ dans
$\mathcal{F}$ défini par
\[
  \mathcal{H}\xrightarrow{v'}\Psi'_*(\mathcal{G})
  \xrightarrow{\Psi'_*(v)}\Psi'_*(\Psi_*(\mathcal{F}))~;
\]
on vérifie que ${v''}^\sharp$ est l'homomorphisme
\[
  \Psi^*({\Psi'}^*(\mathcal{H}))
  \xrightarrow{\Psi^*({v'}^\sharp)}
  \Psi^*(\mathcal{G})
  \xrightarrow{v^\sharp}\mathcal{F}.
\]
\end{env}

\section{Faisceaux quasi-cohérents et faisceaux cohérents}
\label{section:0.5-fr}

\subsection{Faisceaux quasi-cohérents}
\label{subsection:0.5.1-fr}

\begin{env}[5.1.1]
\label{0.5.1.1-fr}
Soient $(X,\mathcal{O}_X)$ un espace annelé, $\mathcal{F}$ un
$\mathcal{O}_X$-Module. La donnée d'un homomorphisme
$u:\mathcal{O}_X\to\mathcal{F}$ de $\mathcal{O}_X$-Modules équivaut à celle
de la section $s=u(1)\in\Gamma(X,\mathcal{F})$. En effet, lorsque $s$ est
donnée, pour toute section $t\in\Gamma(U,\mathcal{O}_X)$, on a nécessairement
$u(t)=t\cdot(s|U)$~; on dit que $u$ est \emph{défini par la section $s$}. Si
maintenant $I$ est un ensemble d'indices quelconque, considérons le faisceau
somme directe $\mathcal{O}_X^{(I)}$, et pour tout $i\in I$, soit $h_i$
l'injection canonique du $i$-ème facteur dans $\mathcal{O}_X^{(I)}$~; on sait
que $u\mapsto(u\circ h_i)$ est un isomorphisme de
$\operatorname{Hom}_{\mathcal{O}_X}(\mathcal{O}_X^{(I)},\mathcal{F})$ sur le
produit
$\bigl(\operatorname{Hom}_{\mathcal{O}_X}(\mathcal{O}_X,\mathcal{F})\bigr)^I$.
Il y a donc correspondance biunivoque canonique entre les homomorphismes
$u:\mathcal{O}_X^{(I)}\to\mathcal{F}$ et les \emph{familles de sections
$(s_i)_{i\in I}$ de $\mathcal{F}$ au-dessus de $X$}. L'homomorphisme $u$
correspondant à $(s_i)$ applique un élément
$(a_i)\in(\Gamma(U,\mathcal{O}_X))^{(I)}$ sur
$\sum_{i\in I}a_i\cdot(s_i|U)$.

On dit que $\mathcal{F}$ est \emph{engendré par la famille $(s_i)$} si
l'homomorphisme $\mathcal{O}_X^{(I)}\to\mathcal{F}$ défini
